\section{Continuous SDM 与 Transformer 组件解释}

经典 SDM 使用 binary address 和 hard radius，Transformer 使用 continuous vector 与 softmax。课程接下来把 SDM 连续化，并问一个更开放的问题：训练好的 Transformer 是否真的落在“好用的 SDM 参数区间”，以及 FFN、LayerNorm 和 memory capacity 能否统一放进这套框架。

\subsection{从 binary radius 到连续 kernel}

Continuous SDM 页说明，不必把输入强制量化为 bit。对归一化实值向量，直接使用指数 dot-product kernel 即可定义连续读出：
$$
\operatorname{CSDM}(q)=
\sum_\mu\operatorname{softmax}_\mu(\beta p_\mu^{\top}q)v_\mu.
$$
这与 attention 形式相同，也可被多层感知机或 associative-memory network 实现。连续化保留“相近地址贡献更大”的原则，却把硬半径变成平滑权重。

\lecturefigure{slide-31-continuous-sdm.jpg}{Continuous SDM：从 binary Hamming space 到实值向量}{官方录像恢复幻灯片 V031；00:21:32}

\begin{knowledgebox}{什么是 effective beta}
$\beta$ 决定 retrieval sharpness：小 $\beta$ 让许多 memory 平均参与，大 $\beta$ 让最相似 memory 占主导。在 Transformer 中，显式 scale、key/query norms 与 LayerNorm 都会共同改变有效温度，所以测量 learned coefficient 应看真实 logits，而不只看公式常数。
\end{knowledgebox}

\subsection{训练后的 Attention 像“好”的 SDM 吗？}

问题页提出反向检验：如果 attention 实现 SDM，训练后参数是否对应容量与噪声鲁棒性较好的半径/temperature？Learned beta 页展示不同 head 的 coefficient 分布，说明模型并没有统一温度。head 可能承担复制、局部语法、长程检索等不同任务，因此最佳 retrieval sharpness 也可能不同。

\lecturefigure{slide-32-transformer-sdm-question.jpg}{开放问题：训练后的 Transformer 是否对应优秀 SDM 参数？}{官方录像恢复幻灯片 V032；00:22:18}

\lecturefigure{slide-33-learned-beta.jpg}{不同 attention heads 学到的 effective beta}{官方录像恢复幻灯片 V033；00:23:44}

\begin{warningbox}{不要把 entropy 当作单调质量指标}
高 entropy 可能表示需要整合多条记忆，也可能表示检索不聚焦；低 entropy 可能是精准召回，也可能是错误过度自信。SDM 视角提供容量/噪声解释，但仍要结合 head 的功能和任务损失判断。
\end{warningbox}

\subsection{FFN、LayerNorm 与 memory framework}

Transformer components 页把 broader architecture 放进 SDM 框架：FFN 可被视为一组长期 key--value memory；LayerNorm 帮助稳定向量范数与 coefficient；attention 是显式 pattern-based retrieval。最后 summary 页把结论与未来问题并列，强调这些对应关系能够产生研究假设，却不是一套已经证完的统一理论。

\lecturefigure{slide-34-transformer-components-converge.jpg}{Transformer 组件如何汇入 SDM framework}{官方录像恢复幻灯片 V034；00:25:06}

\lecturefigure{slide-35-sdm-summary.jpg}{第一讲理论部分总结与开放问题}{官方录像恢复幻灯片 V035；00:26:20}

\begin{longtable}{@{}L{0.20\textwidth}L{0.31\textwidth}L{0.37\textwidth}@{}}
\toprule
Transformer 组件 & SDM/记忆视角 & 必须保留的限制 \\
\midrule
Attention & 当前上下文中的精确 pattern-based retrieval & 受 context window 约束，存储所有 keys/values \\
FFN & 参数中叠加的长期 key--value associations & 检索更噪、更难逐条追溯，不等于完整数据集存储 \\
LayerNorm & 稳定向量 norm，使距离/点积对应更可控 & 实际 norms 与 learned scales 仍会改变 temperature \\
Multi-head & 多组不同表示子空间与 retrieval sharpness & 不能直接等同于脑区 microzones \\
\bottomrule
\end{longtable}

\teachervoice{问答中，讲者区分 neuron-centric 与 pattern-centric SDM：attention 保留所有 pattern location，因此相似度精确但有 context window；长期分布式记忆只保留 noisy superposition，容量更大却难以精确三角定位原 pattern。这个 tradeoff 比旧稿虚构的生产监控指标更接近课堂真正的工程启示。}

\subsection{本章小结}

连续 SDM 将硬半径变成指数 kernel，与 attention 共享形式。Learned beta、LayerNorm 和 FFN 解释把框架扩展到完整 Transformer，但每个对应都应作为可检验假设。Pattern memory 的精确短上下文与 distributed memory 的高容量噪声，是课程提出的核心权衡。

